\question[4]
Gegeben ist der folgende Algorithmus in Pseudocode.
\begin{lstlisting}[language={}]
n = 13
s = 0
solange n > 0:
    s = s + n % 2
    n = n // 2
gib s aus
\end{lstlisting}
\begin{parts}
\part Erstelle ein Ablaufprotokoll für n und s. Was wird ausgegeben?
\part Was berechnet der Algorithmus für eine beliebige positive ganze Zahl n?
\end{parts}
\begin{solutionbox}{3.5cm}
\begin{tabular}{l|ccccc}
Durchlauf & Start & 1 & 2 & 3 & 4 \\
\hline
n & 13 & 6 & 3 & 1 & 0 \\
s & 0 & 1 & 1 & 2 & 3 \\
\end{tabular}

Ausgabe: 3. Der Algorithmus zählt die Einsen in der Dualdarstellung von n.
\end{solutionbox}
